% Einheit 3 von 4 – Grundwert berechnen (Katalog Einheit 4), Kl. 7, schwach
\clearpage
\hypertarget{e3}{}
\einheitenkopf*{Einheit 3 von 4 · Grundwert berechnen}
\zweigzeile{Hier lernst du, aus einem Teil und seiner Prozentzahl auszurechnen, wie groß das Ganze ist · neu in diesem Jahr · P10 oft · baut auf: Einheit 1 und 2 (für die gemischten Aufgaben), Dreisatz (Kennst du schon)}
\setcounter{aufgabe}{30}

\begin{aufgabe}{Ich erkenne, was gegeben und was gesucht ist. Das Ganze heißt Grundwert, der Teil heißt Prozentwert, die Prozentzahl heißt Prozentsatz. Rechne nicht. Kreuze an, was gesucht ist.}
\swb{\small gefärbt: Prozentwert \\ ganzer Streifen: Grundwert}{\streifen{30}{$0\,\%$}{$100\,\%$}}
\begin{teile}
\teil Wie viel sind 20\,\% von 50\,€? \\ \kreuz{Prozentwert}\quad\kreuz{Prozentsatz}\quad\kreuz{Grundwert}
\teil Wie viel Prozent sind 9 von 36? \\ \kreuz{Prozentwert}\quad\kreuz{Prozentsatz}\quad\kreuz{Grundwert}
\teil 12\,€ sind 25\,\%. Wie viel sind 100\,\%? \\ \kreuz{Prozentwert}\quad\kreuz{Prozentsatz}\quad\kreuz{Grundwert}
\teil Von 50 Äpfeln sind 6 faul. Wie viel Prozent der Äpfel sind faul? \\ \kreuz{Prozentwert}\quad\kreuz{Prozentsatz}\quad\kreuz{Grundwert}
\teil 14 Kinder sind 70\,\% der Gruppe. Wie viele Kinder hat die Gruppe? \\ \kreuz{Prozentwert}\quad\kreuz{Prozentsatz}\quad\kreuz{Grundwert}
\end{teile}
\end{aufgabe}

\begin{aufgabe}{Ich kann am Streifen vom Teil auf das Ganze schließen. Der gefärbte Teil ist gegeben. Rechne aus, wie viel der ganze Streifen ist.}
\swb{\small 30\,\% sind 18\,kg.\\[2pt]\beispiel{1 \text{ Kästchen} &= 18\,\text{kg} : 3 = 6\,\text{kg} \\ 10 \text{ Kästchen} &= 6\,\text{kg} \cdot 10 = 60\,\text{kg}}}{\streifen{30}{$0$\,kg}{$?$}}
\swz{10\,\% sind 5\,€.}{\streifen{10}{$0\,€$}{$?$}}
\swz{20\,\% sind 14\,m.}{\streifen{20}{$0$\,m}{$?$}}
\swz{50\,\% sind 13\,kg.}{\streifen{50}{$0$\,kg}{$?$}}
\swz{25\,\% sind 11\,€.}{\streifen[4]{25}{$0\,€$}{$?$}}
\end{aufgabe}

\begin{aufgabe}{Ich kann aus einem Teil das Ganze ausrechnen. Überlege, wie oft die Prozentzahl in 100\,\% passt. So oft nimmst du den Teil mal. Das Ergebnis ist der Grundwert.}
\swb{\small 50\,\% sind 21\,€.\\[2pt]\beispiel{50\,\% \cdot 2 &= 100\,\% \\ 21\,€ \cdot 2 &= 42\,€}}{\streifen{50}{$0\,€$}{$?$}}
\anw{Der Teil ist jetzt immer 9\,€.}
\swz{9\,€ sind 50\,\%.}{\streifen{50}{$0\,€$}{$?$}}
\swz{9\,€ sind 25\,\%.}{\streifen[20]{25}{$0\,€$}{$?$}}
\swz{9\,€ sind 20\,\%.}{\streifen{20}{$0\,€$}{$?$}}
\swz{9\,€ sind 10\,\%.}{\streifen{10}{$0\,€$}{$?$}}
\swfrage{Was bleibt gleich, was ändert sich?}
\end{aufgabe}

\begin{aufgabe}{Ich kann aus einem Teil das Ganze ausrechnen – weiter. Rechne zuerst 1\,\% aus: Teil : Prozentzahl. Nimm dann mal 100.}
\swb{\small 3\,\% sind 12\,kg.\\[2pt]\beispiel{1\,\% &= 12\,\text{kg} : 3 = 4\,\text{kg} \\ 100\,\% &= 4\,\text{kg} \cdot 100 = 400\,\text{kg}}}{\streifen[0]{3}{$0$\,kg}{$?$}}
\swz{6\,\% sind 18\,€.}{\streifen[0]{6}{$0\,€$}{$?$}}
\anw{Rechne mit dem Taschenrechner.}
\swz{35\,\% sind 91\,kg.}{\streifen[20]{35}{$0$\,kg}{$?$}}
\anw{Lies genau: Was ist der Teil, was ist die Prozentzahl?}
\swz{In der Klasse 7b fahren 18 Kinder mit dem Bus. Das sind 75\,\% der Klasse. Wie viele Kinder sind in der Klasse?}{\streifen[4]{75}{$0$}{alle Kinder}}
\end{aufgabe}

\begin{aufgabe}{Ich kann unterscheiden, was gesucht ist, und dann richtig rechnen. Kreuze zuerst an, was gesucht ist. Rechne dann.}
\swz{Wie viel sind 30\,\% von 80\,kg?}{\kreuz{Prozentwert}\quad\kreuz{Prozentsatz}\quad\kreuz{Grundwert}\par\smallskip\streifenleer}
\swz{Ein Buch hat 250 Seiten. Emma hat 75 Seiten gelesen. Wie viel Prozent hat sie gelesen?}{\kreuz{Prozentwert}\quad\kreuz{Prozentsatz}\quad\kreuz{Grundwert}\par\smallskip\streifenleer}
\swz{In einer Dose sind noch 24 Kekse. Das sind 40\,\% der Kekse vom Anfang. Wie viele Kekse waren am Anfang in der Dose?}{\kreuz{Prozentwert}\quad\kreuz{Prozentsatz}\quad\kreuz{Grundwert}\par\smallskip\streifenleer}
\begin{teile}
\teil Beim Kauf eines Rucksacks spart Jonas 12\,€. Das sind 15\,\% des alten Preises. Wie viel kostete der Rucksack vorher? (P10 2025)
\end{teile}
\end{aufgabe}

\begin{aufgabe}{Ich finde den Fehler, wenn das Ganze gesucht ist. 12\,€ sind 25\,\% eines Preises. Gesucht ist der ganze Preis. Sara rechnet so:}
\rechnung{25\,\% \text{ von } 12\,€ &= 3\,€ \\ \text{Der ganze Preis ist } &3\,€.}
\swz{Finde den Fehler, benenne ihn und korrigiere die Rechnung.}{\streifen[4]{25}{$0\,€$}{$?$}}
\swz{7\,kg sind 20\,\%. Wie viel ist das Ganze?}{\streifen{20}{$0$\,kg}{$?$}}
\end{aufgabe}

\begin{aufgabe}{Ich kann begründen, warum ich am Ende mal 100 nehme.}
\swz{Warum nimmt man beim Weg über 1\,\% am Ende mal 100? Begründe.}{\streifen[0]{1}{$0$}{$100\,\%$}}
\swz{Tom sagt: „Wenn 150\,\% eines Betrags 30\,€ sind, dann ist der ganze Betrag kleiner als 30\,€.“ Stimmt das? Begründe.}{\zahlenstrahl[xmin=0,xmax=1.55,xstep=0.1,karo=0.6]{1/100\,\%, 1.5/150\,\%}}
\end{aufgabe}

\begin{aufgabe}{Ich kann ausrechnen, wie viel Geld noch fehlt. Mia spart auf ein Skateboard. Sie hat schon 42\,€ gespart. Das sind 70\,\% des Preises.}
\swz{Wie viel kostet das Skateboard?}{\streifen{70}{$0\,€$}{$?$}}
\swz{Wie viel Euro fehlen Mia noch? Das ist der weiße Teil des Streifens.}{\streifen{70}{$0\,€$}{$60\,€$}}
\swz{Mia bekommt jede Woche 4\,€ Taschengeld und spart alles. Hat sie nach 4 Wochen genug Geld? Begründe.}{}
\end{aufgabe}

\uebersichtskasten{Grundwert: erst 1\,\% (Prozentwert : p), dann mal 100. \\ 36\,€ sind 45\,\%: \ 1\,\% $= 0{,}80$\,€ $\to$ 100\,\% $= 80$\,€ \qquad Streifen: 45\,\% $\mathrel{\widehat{=}}$ 36\,€, 100\,\% $\mathrel{\widehat{=}}$ ? \\ Erst zuordnen: Was ist gegeben (W, p, G), was gesucht? Dann rechnen.}
